真空中 \(z=0\) 平面放置一张无限薄、均匀且局域响应的二维导电层，法向取 \(+\hat{\mathbf z}\)。单色平面波以 \(e^{i(kz-\omega t)}\) 约定从 \(z<0\) 正入射，电场切向复振幅为 \(\mathbf E_i=E_0\hat{\mathbf x}\)（\(E_0>0\)）。两侧介质最初均为真空，阻抗 \(Z_0=\sqrt{\mu_0/\epsilon_0}\)。忽略层厚、空间色散、外加电流源及磁化面电流。二维欧姆定律为
\[
\mathbf K=\boldsymbol{\sigma}\mathbf E_t,\qquad
\boldsymbol{\sigma}=\sigma_d\mathbb I+\sigma_H\mathsf J,\qquad
\mathsf J=\begin{pmatrix}0&-1\\1&0\end{pmatrix},\quad
\sigma_d\ge0,\quad\sigma_H\in\mathbb R .
\]
这里 \(\mathbf E_t\) 同时是薄层处的连续切向电场，\(\sigma_d,\sigma_H\) 均取实数；令 \(a=Z_0\sigma_d/2\)、\(b=Z_0\sigma_H/2\)。以下功率比均指时间平均、单位面积的法向能流比。
(1) 严格使用切向 \(\mathbf E\) 连续及 \(\hat{\mathbf z}\times(\mathbf H_{z>0}-\mathbf H_{z<0})=\mathbf K\) 推导反射、透射的 \(2\times2\) Jones 矩阵 \(\mathcal R,\mathcal T\)。求两圆偏振本征通道的振幅本征值，并说明其复数相位的符号约定。
(2) 令 \(\sigma_d=0\) 而 \(\sigma_H\ne0\)。求线偏振入射后的透射和反射电场、两者偏振方向及功率 \(T,R\)；证明反射与透射的线偏振方向正交，核对能量守恒并给出 \(b\to0\) 极限。
(3) 恢复 \(\sigma_d\ge0\)。分别求 \(T,R,A\)（\(A\) 为层内焦耳耗散份额），证明 \(R+T+A=1\) 与 \(A\ge0\)。固定霍尔参数 \(b\) 时，在 \(a\ge0\) 上严格求使 \(A\) 最大的纵向电导及最大吸收率，并讨论 \(b=0\) 和 \(a\to\infty\) 极限。
(4) 改为层两侧分别是折射率 \(n_1,n_2>0\) 的无耗、非磁性均匀介质，波仍从介质1正入射，且令 \(\sigma_d=0\)。从边界条件重新求 \(\mathcal T_{12},\mathcal R_{12}\)，推导偏振旋转角及能流比 \(T_{12},R_{12}\)。若在低频无耗整数霍尔平台 \(\sigma_H=\nu e^2/h_{\rm P}\)（\(\nu\) 为整数、\(h_{\rm P}=2\pi\hbar\)），用精细结构常数 \(\alpha_{\rm fs}=e^2/(4\pi\epsilon_0\hbar c)\) 表示旋转角，核对无导电层的 Fresnel 极限及真空小 \(|\nu\alpha_{\rm fs}|\) 极限。
